\section{Introduction}

\subsection{Use case}
We chose Use Case 1, Ride-Hailing Fleet Operations. Our operator, FleetPulse, runs 30 vehicles in
Colombo: 13 tuk-tuks, 10 mini cars and 7 cars. The operator owns the vehicles. It pays for fuel and
maintenance and gives each driver 30\% of the fares. We divided the city into 12 zones, such as
Fort-Pettah, Kollupitiya, Borella, Wellawatte and Nugegoda. The use case asks one question with two parts:

\begin{quote}\itshape
``What is fleet utilization and earnings by area/time-of-day right now, and which vehicles are becoming
unprofitable once yesterday's fuel/maintenance costs are factored in?''
\end{quote}

The first part needs fresh numbers within seconds. The second part needs exact money numbers, but only once
a day. This difference is the main reason for our design.

\subsection{Business requirements}
Table~\ref{tab:requirements} lists the requirements we took from the use case. The last two columns show how
fast and how exact each answer must be.

\begin{table}[!htb]
\centering\small
\caption{Business requirements taken from the use case.}
\label{tab:requirements}
\begin{tabularx}{\textwidth}{@{}l L P{2.9cm} P{3.6cm}@{}}
\toprule
\textbf{ID} & \textbf{Requirement} & \textbf{How fast} & \textbf{How exact} \\
\midrule
R1 & Live fleet status: vehicles online, on a trip and idle & a few seconds & approximate is fine \\
R2 & Utilisation and fares by zone and hour of day & seconds & small errors are fine \\
R3 & Alert when a vehicle stays idle for a long time & simulated minutes & one alert per idle period \\
R4 & Daily profit of each vehicle after fuel and maintenance & once per day & exact, it is money \\
R5 & Find vehicles that are becoming unprofitable & once per day & exact, needs history \\
R6 & Recompute a day when late data or a corrected file arrives & on request & exact and repeatable \\
R7 & Show the health of every pipeline stage & seconds & logs, metrics and alerts \\
R8 & Run on one laptop with 8\,GB of memory & always & low memory and CPU \\
\bottomrule
\end{tabularx}
\end{table}

\subsection{Simulated clock and assumptions}
A real business day is long. To show many days in a short demo, we compress time. One real second is two
simulated minutes (a speed-up of 120), so 30 real seconds are one simulated hour and 12 real minutes are one
simulated day. The simulation starts on Tuesday 1~September~2026 at 06:00, so day~1 is 18 hours long and
later days are 24 hours long. All simulated times are Sri Lanka local time. The clock is stored in one
database row, so every container agrees on the simulated ``now''. If the whole system stops for more than
five minutes, simulated time pauses and continues after the restart.

We also made these assumptions:
\begin{itemize}
\item Both data sources are Python scripts. The numbers are realistic but they are not real company data.
\item Money is in Sri Lankan rupees (LKR). Petrol costs about LKR 305 per litre.
\item The daily cost file arrives 20 to 90 simulated minutes after midnight (10 to 45 real seconds).
\item Kafka, Spark and PostgreSQL each run as a single node. This is enough for about 30 events per second.
\end{itemize}
